\documentclass[11pt]{amsart}
\usepackage[utf8]{inputenc}
\usepackage[T1]{fontenc}
\usepackage{lmodern}
\usepackage{amsmath,amssymb,amsthm}
\usepackage{mathtools}
\usepackage{booktabs}
\usepackage{array}
\usepackage{tabularx}
\usepackage{enumitem}
\usepackage[margin=1.1in]{geometry}
\usepackage[hidelinks]{hyperref}
\hypersetup{
  pdftitle={Extremal coefficients of Pell representation polynomials and integer fibres in the silver ratio},
  pdfauthor={HYDRA Research Project},
  pdfsubject={HGS-SILVER full manuscript, pre-submission draft},
  pdfkeywords={Pell numbers; silver ratio; representation function; extremal coefficients; numeration systems}
}

\newtheorem{lemma}{Lemma}[section]
\newtheorem{proposition}[lemma]{Proposition}
\newtheorem{corollary}[lemma]{Corollary}
\theoremstyle{definition}
\newtheorem{definition}[lemma]{Definition}
\theoremstyle{remark}
\newtheorem{remark}[lemma]{Remark}
\theoremstyle{plain}
\newtheorem*{thmPETone}{Theorem PET-1}
\newtheorem*{thmPETtwo}{Theorem PET-2}
\newtheorem*{thmPETthree}{Theorem PET-3}
\newtheorem*{thmPETfour}{Theorem PET-4}
\newtheorem*{thmSTTone}{Theorem STT-1}
\newtheorem*{thmSTTtwo}{Theorem STT-2}
\newtheorem*{thmSTTthree}{Theorem STT-3}

\numberwithin{equation}{section}
\raggedbottom

\newcommand{\Z}{\mathbb{Z}}
\newcommand{\Q}{\mathbb{Q}}
\newcommand{\floor}[1]{\lfloor #1\rfloor}
\newcommand{\argmax}{\operatorname*{argmax}}
\newcommand{\CC}{\mathsf{C}}
\newcommand{\OL}{\mathsf{OL}}
\newcommand{\OR}{\mathsf{OR}}
\newcommand{\IL}{\mathsf{IL}}
\newcommand{\IR}{\mathsf{IR}}
\newcommand{\Fam}{\mathcal{F}}
\newcommand{\PETone}{\hyperref[thm:pet1]{Theorem~PET-1}}
\newcommand{\PETtwo}{\hyperref[thm:pet2]{Theorem~PET-2}}
\newcommand{\PETthree}{\hyperref[thm:pet3]{Theorem~PET-3}}
\newcommand{\PETfour}{\hyperref[thm:pet4]{Theorem~PET-4}}
\newcommand{\STTone}{\hyperref[thm:stt1]{Theorem~STT-1}}
\newcommand{\STTtwo}{\hyperref[thm:stt2]{Theorem~STT-2}}
\newcommand{\STTthree}{\hyperref[thm:stt3]{Theorem~STT-3}}

\title[Pell representation polynomials and silver-ratio fibres]{Extremal coefficients of Pell representation\\ polynomials and integer fibres in the silver ratio}

\author[HYDRA Research Project]{HYDRA Research Project}
\thanks{Human-directed, AI-authored mathematical research. The proofs have not undergone independent human mathematical verification.}

\subjclass[2020]{Primary 11B39; Secondary 11A63, 05A15, 05A17}
\keywords{Pell numbers; silver ratio; representation function; restricted-digit expansions;
extremal coefficients; numeration systems}

\begin{document}

\begin{abstract}
Let $P_n$ be the Pell numbers and, for $\ell\ge1$, let $\rho_\ell(n)$ be the coefficient of $x^n$
in $G_\ell(x)=\prod_{j=2}^{\ell+1}(1+x^{P_j}+x^{2P_j})$, that is, the number of words
$(d_2,\dots,d_{\ell+1})\in\{0,1,2\}^\ell$ with $\sum_j d_jP_j=n$. We prove that
$B_\ell=\max_n\rho_\ell(n)=2^{\lfloor(\ell-1)/2\rfloor}$ for every $\ell\ge1$, and we determine the
set $M_\ell$ of all maximizers for every length. For odd $\ell\ge3$ it consists of the four integers
$u_\ell,\ u_\ell+2,\ u_\ell+P_{\ell+1},\ u_\ell+P_{\ell+1}+2$ with $u_\ell=(P_{\ell+2}-5)/2$, and
$M_1=\{0,2,4\}$. For even lengths $|M_2|=9$ and $|M_{2k}|=16k-8$ for $k\ge2$; moreover $M_{2k}$ is
the union of two explicit six-element cores and four explicit chains of two-element blocks. All
proofs rest on one structural fact: when two consecutive Pell weights are adjoined, the nine shifted
copies of the previous coefficient sequence overlap at most two at a time. As an application, let
$\delta=1+\sqrt2$ and let $\nu_m(t)$ be the number of strings
$(d_k)_{k=-2m}^{2m}\in\{0,1,2\}^{4m+1}$ with $\sum_kd_k\delta^k=t$. We prove the exact identity
$\nu_m(t)=\rho_{4m+1}(tP_{2m+2})$ for every integer $t$. Consequently
$\max_{t\in\mathbb Z}\nu_m(t)=4^m$, attained only at $t=\sum_{j=-m}^{m}\delta^{2j}=P_{2m+1}+P_{2m}$
when $m\ge1$ and at $t\in\{0,1,2\}$ when $m=0$, and every fibre over $\mathbb Z[\sqrt2]$ in the same
window contains at most $4^m$ strings. The results and proofs constitute human-directed, AI-authored mathematical research
and have been checked only by AI-model audits; no human review has taken place, and
novelty is not established.
\end{abstract}

\maketitle

%======================================================================
\section{Introduction and related work}\label{sec:intro}
%======================================================================

Counting the representations of an integer in a redundant numeration system is a classical problem.
For sums of distinct Fibonacci numbers the representation function $R(n)$ has been studied in
detail. Koc\'abov\'a, Mas\'akov\'a and Pelantov\'a~\cite{KMP2005} determine the maximum of $R(n)$
between consecutive Fibonacci numbers and describe the integers at which it is attained, and
Stockmeyer~\cite{Stockmeyer2009} proves $R(n)\le\sqrt{n+1}$, with equality exactly at
$n=F_m^2-1$. The same programme---the maximum of the representation function at a fixed numeration
length, together with a description of its maximizers---is carried out for the $m$-bonacci systems
in~\cite{KMP2007}. For expansions in Pisot bases, Grabner and Steiner~\cite{GrabnerSteiner2011}
study expansions of minimal weight; in the Fibonacci case with digits $\{-1,0,1\}$ they obtain the
exact maximum $2^{\floor{(k-1)/3}}$ for the number of such expansions of length
$k$~\cite[Theorem~4.6]{GrabnerSteiner2011}. Edson~\cite{Edson2013} develops methods for computing
numbers of representations in linear numeration systems.

For the Pell numbers $P_0=0$, $P_1=1$, $P_{n+2}=2P_{n+1}+P_n$, the literature we consulted
concentrates on \emph{distinguished single representations}. Horadam~\cite{Horadam1994} describes the
\emph{maximal} (lazy) representation, in which every positive integer is written with the largest
possible number of Pell summands; the \emph{minimal} (greedy) representation and the recursive
counting of partitions into Pell parts of \emph{unbounded} multiplicity are treated by Knopfmacher
and Robbins~\cite{KnopfmacherRobbins2004}; both normalizations appear in~\cite{Horadam1996}. They are
recalled in Section~\ref{sec:terminology} to separate them from the quantity studied here.

In this paper we study the bounded Pell representation polynomials
\[
G_\ell(x)=\prod_{j=2}^{\ell+1}\bigl(1+x^{P_j}+x^{2P_j}\bigr),\qquad \rho_\ell(n)=[x^n]\,G_\ell(x),
\qquad \ell\ge1,
\]
so that $\rho_\ell(n)$ is the number of words $(d_2,\dots,d_{\ell+1})\in\{0,1,2\}^\ell$ with
$\sum_jd_jP_j=n$. We ask for the maximum coefficient $B_\ell=\max_n\rho_\ell(n)$ and for the set
$M_\ell=\{n:\rho_\ell(n)=B_\ell\}$ of all maximizers.

\subsection{Main results}
Our motivation is a question about the silver ratio $\delta=1+\sqrt2$: how many ternary digit strings
of length $4m+1$, centred at the exponent~$0$, can represent the same integer? Let
\[
\nu_m(t)=\#\Bigl\{(d_k)_{k=-2m}^{2m}\in\{0,1,2\}^{4m+1}:\ \sum_kd_k\delta^k=t\Bigr\},\qquad
N_m=\max_{t\in\Z}\nu_m(t).
\]
The results are stated under seven canonical identifiers, which we keep throughout.
\begin{enumerate}[label=\textbullet,leftmargin=1.6em]
\item \PETone{} (Section~\ref{sec:pet1}): $B_\ell=2^{\floor{(\ell-1)/2}}$ for every $\ell\ge1$.
\item \PETtwo{} (Section~\ref{sec:pet2}): for every odd $\ell\ge1$, the complete list of maximizers;
four integers for $\ell\ge3$, and $M_1=\{0,2,4\}$.
\item \PETthree{} (Section~\ref{sec:pet3}): $|M_2|=9$ and $|M_{2k}|=16k-8$ for every $k\ge2$; the
case $k=1$ is exceptional.
\item \PETfour{} (Section~\ref{sec:pet4}): an explicit description of $M_{2k}$ for every $k\ge1$ as
the union of two cores and four chains of two-element blocks; the families are pairwise disjoint for
$k\ge2$, while for $k=1$ the two cores meet in $\{5,7,9\}$.
\item \STTone{} (Section~\ref{sec:stt1}): $\nu_m(t)=\rho_{4m+1}(tP_{2m+2})$ for all $m\ge0$ and
$t\in\Z$, by an explicit bijection of words.
\item \STTtwo{} (Section~\ref{sec:stt2}): $N_m=4^m$ for every $m\ge0$; the maximum is attained only
at the integer $C_m=\sum_{j=-m}^m\delta^{2j}$ when $m\ge1$, and at $t\in\{0,1,2\}$ when $m=0$.
\item \STTthree{} (Section~\ref{sec:stt3}): every fibre over $\Z[\sqrt2]$ in the window
$\{-2m,\dots,2m\}$ contains at most $4^m$ strings.
\end{enumerate}
\STTone{} reduces the silver-ratio question exactly to bounded Pell representations, and
\PETone{} and \PETtwo{} solve that problem. The even-length results \PETthree{} and \PETfour{} are not
needed for the silver-ratio application; they complete the description of $M_\ell$.

All proofs rest on one structural fact (Lemma~\ref{lem:window}): when the two weights $P_{\ell+2}$
and $P_{\ell+3}$ are adjoined to $G_\ell$, the nine shifted copies of the coefficient sequence of
$G_\ell$ overlap at most two at a time, and only when they are adjacent. This forces the doubling
$B_{\ell+2}=2B_\ell$ and turns the maximizer sets into an exact set recursion.

\subsection{Scope and status}
The statements concern the Pell recurrence, the digit set $\{0,1,2\}$, the weight window
$P_2,\dots,P_{\ell+1}$, the exponent window $\{-2m,\dots,2m\}$ and the target domains stated in each
theorem; they do not extend to other recurrences, alphabets, windows or quadratic units. We do not
claim priority: the statements were not found in the literature we consulted, but that search had
documented gaps, which are listed in Section~\ref{sec:prior}. The mathematical work is human-directed, AI-authored research
and has been checked by AI-model audits only; the statement at the
end of the paper describes which system did what. No human mathematician has reviewed the proofs.

\subsection{Organization}
Section~\ref{sec:defs} fixes notation and collects the elementary identities used later.
Sections~\ref{sec:pet1}--\ref{sec:pet4} prove the four Pell theorems, and
Sections~\ref{sec:stt1}--\ref{sec:stt3} the three silver-ratio theorems. Section~\ref{sec:discussion}
discusses scope, limitations, the literature and the verification status. Appendix~\ref{app:examples}
lists small cases, Appendix~\ref{app:status} summarizes what is and is not established, and
Appendix~\ref{app:aliases} records the historical labels of the theorems and the role of computer
checks.

%======================================================================
\section{Definitions and Pell representation polynomials}\label{sec:defs}
%======================================================================

\subsection{Pell numbers}
Let $P_0=0$, $P_1=1$ and $P_{n+2}=2P_{n+1}+P_n$; thus
$P_2=2$, $P_3=5$, $P_4=12$, $P_5=29$, $P_6=70$, $P_7=169$. We extend the sequence to all $n\in\Z$ by
$P_{-n}=(-1)^{n+1}P_n$. The recurrence then holds for all $n\in\Z$: for $n=1$ it reads
$P_1-2P_0=1=P_{-1}$, and for $n\ge2$ one has
$P_{-n+2}-2P_{-n+1}=(-1)^{n-1}P_{n-2}+2(-1)^{n-1}P_{n-1}=(-1)^{n+1}(P_{n-2}+2P_{n-1})=P_{-n}$.

\begin{lemma}[sums and parity]\label{lem:sums}
\begin{enumerate}[label=\textup{(\alph*)}]
\item $\sum_{j=1}^{r}P_j=(P_{r+1}+P_r-1)/2$ for every $r\ge0$.
\item $P_n\equiv n\pmod 2$ for every $n\ge0$.
\item $2P_n=P_{n+1}-P_{n-1}$ for every $n\ge1$. Consequently, for every $m\ge0$,
\[
2\sum_{i=1}^{m}P_{2i}=P_{2m+1}-1,\qquad 2\sum_{i=1}^{m}P_{2i-1}=P_{2m},\qquad
2\sum_{i=0}^{m}P_{2i+1}=P_{2m+2}.
\]
\end{enumerate}
\end{lemma}

\begin{proof}
(a) For $r=0$ both sides vanish. Passing from $r-1$ to $r$, the right-hand side increases by
$(P_{r+1}-P_{r-1})/2=P_r$, which is the $r$-th summand.
(b) $P_0$ and $P_1$ have the stated parities, and $P_{n+2}=2P_{n+1}+P_n\equiv P_n\pmod2$.
(c) $P_{n+1}-P_{n-1}=2P_n$ is the recurrence. The three sums telescope:
$\sum_{i=1}^m(P_{2i+1}-P_{2i-1})=P_{2m+1}-P_1$, $\sum_{i=1}^m(P_{2i}-P_{2i-2})=P_{2m}-P_0$ and
$\sum_{i=0}^m(P_{2i+2}-P_{2i})=P_{2m+2}-P_0$.
\end{proof}

\begin{lemma}[addition and Cassini identities]\label{lem:cassini}
For integers $a,b\ge1$, $P_{a+b}=P_aP_{b+1}+P_{a-1}P_b$. For $n\ge1$,
$P_n^2-P_{n-1}P_{n+1}=(-1)^{n+1}$.
\end{lemma}

\begin{proof}
Induction on $n$ gives
\[
Q=\begin{pmatrix}2&1\\1&0\end{pmatrix},\qquad
Q^n=\begin{pmatrix}P_{n+1}&P_n\\P_n&P_{n-1}\end{pmatrix}\quad(n\ge1).
\]
The lower-left entry of $Q^aQ^b=Q^{a+b}$ is the addition identity. Taking determinants,
$P_{n+1}P_{n-1}-P_n^2=(\det Q)^n=(-1)^n$.
\end{proof}

\begin{lemma}[powers of the silver ratio]\label{lem:powers}
Let $\delta=1+\sqrt2$ and $\lambda=1-\sqrt2=-\delta^{-1}$. For every $k\in\Z$,
\[
\delta^k=P_k\,\delta+P_{k-1},\qquad \lambda^k=P_k\,\lambda+P_{k-1}.
\]
Consequently $P_k=(\delta^k-\lambda^k)/(2\sqrt2)$ and $P_k+P_{k-1}=(\delta^k+\lambda^k)/2$.
\end{lemma}

\begin{proof}
Both $\delta$ and $\lambda$ are roots of $x^2=2x+1$. The formulas hold for $k=0$ (as $P_{-1}=1$) and
$k=1$, and both sides satisfy the Pell recurrence in $k$, forwards and backwards. Subtracting the two
formulas gives $\delta^k-\lambda^k=P_k(\delta-\lambda)=2\sqrt2\,P_k$; adding them gives
$\delta^k+\lambda^k=P_k(\delta+\lambda)+2P_{k-1}=2P_k+2P_{k-1}$.
\end{proof}

\subsection{Two distinguished single representations}\label{sec:terminology}
The word ``canonical'' is used in more than one sense for Pell representations, and we avoid it for
representations. Two normalizations are in use.
\begin{itemize}[leftmargin=1.6em]
\item \emph{Greedy (minimal) representation} \cite{KnopfmacherRobbins2004,Horadam1996}: every positive
integer has a unique representation $N=\sum_i\alpha_iP_i$ with $\alpha_i\in\{0,1,2\}$, subject to
$\alpha_i=2\Rightarrow i\ge2$ and $\alpha_{i-1}=0$.
\item \emph{Lazy (maximal) representation} \cite{Horadam1994,Horadam1996}: every positive integer has
a unique representation $N=\sum_i\beta_iP_i$ with $\beta_i\in\{0,1,2\}$, finitely many non-zero,
subject to $\beta_i=0\Rightarrow\beta_{i-1}=2$, and $\beta_k\in\{1,2\}$ for the largest index $k$.
\end{itemize}
The ``maximal representation'' of~\cite{Horadam1994} maximizes the number of summands of a
\emph{single} representation. Every statement in this paper concerns a different quantity: the
\emph{number} of representations of a fixed integer, that is, the size of a whole fibre.

\subsection{The Pell counting model}
For $\ell\ge1$ let
\begin{equation}\label{eq:model}
G_\ell(x)=\prod_{j=2}^{\ell+1}\bigl(1+x^{P_j}+x^{2P_j}\bigr),\qquad
\rho_\ell(n)=[x^n]\,G_\ell(x),
\end{equation}
\[
B_\ell=\max_{n\in\Z}\rho_\ell(n),\qquad M_\ell=\{n:\rho_\ell(n)=B_\ell\},
\]
and let $D_\ell=\deg G_\ell=2\sum_{j=2}^{\ell+1}P_j$. We set $\rho_\ell(n)=0$ for $n\notin[0,D_\ell]$.
Equivalently, $\rho_\ell(n)$ is the number of words $d\in\{0,1,2\}^{\{2,\dots,\ell+1\}}$ with
$\sum_jd_jP_j=n$. The weight $P_1=1$ is never used: the window always begins at $P_2=2$.

\begin{lemma}[degree and symmetry]\label{lem:degree}
For every $\ell\ge1$, $D_\ell=P_{\ell+2}+P_{\ell+1}-3$, and $\rho_\ell(n)=\rho_\ell(D_\ell-n)$ for
all $n\in\Z$. In particular $n\in M_\ell$ if and only if $D_\ell-n\in M_\ell$.
\end{lemma}

\begin{proof}
By Lemma~\ref{lem:sums}(a),
$D_\ell=2\bigl(\sum_{j=1}^{\ell+1}P_j-P_1\bigr)=(P_{\ell+2}+P_{\ell+1}-1)-2$. The map
$d_j\mapsto2-d_j$ is a bijection of $\{0,1,2\}$ and carries a word of value $n$ to a word of value
$\sum_j(2-d_j)P_j=D_\ell-n$.
\end{proof}

\subsection{The silver-ratio model}\label{sec:silvermodel}
For $m\ge0$ put $s=2m+2$ and
\[
R_m=\rho_{4m+1},\qquad H_m=D_{4m+1}=P_{4m+3}+P_{4m+2}-3,\qquad C_m=\sum_{j=-m}^{m}\delta^{2j},
\]
and let $\nu_m$ and $N_m$ be as in Section~\ref{sec:intro}. This is a \emph{ternary, symmetric-window}
model: the digits are $0,1,2$, the exponent window is $\{-2m,\dots,2m\}$, and in $\nu_m(t)$ the
target $t$ is a rational integer. Since every $\delta^k$ lies in $\Z[\sqrt2]$, every value
$\sum_kd_k\delta^k$ lies in $\Z[\sqrt2]$. The integer $s=2m+2$ is always even; this parity is used in
the proofs of \STTone{} and of Lemma~\ref{lem:residue}.

%======================================================================
\section{Exact maximum coefficient}\label{sec:pet1}
%======================================================================

Fix $\ell\ge1$ and write
\[
U=P_{\ell+1},\qquad V=P_{\ell+2},\qquad W=P_{\ell+3}=2V+U,\qquad D=D_\ell=U+V-3,
\]
the last identity by Lemma~\ref{lem:degree}. Since $V=2U+P_\ell$, we have $V-U=U+P_\ell\ge3$, with
equality only for $\ell=1$; in particular $V>U+2$. Since
$G_{\ell+2}=G_\ell\,(1+x^V+x^{2V})(1+x^W+x^{2W})$,
\begin{equation}\label{eq:twostep}
\rho_{\ell+2}(n)=\sum_{a,b\in\{0,1,2\}}\rho_\ell(n-aV-bW)\qquad(n\in\Z).
\end{equation}
This is the coefficient identity for a product of polynomials: distinct pairs $(a,b)$ correspond to
distinct words, so no word is lost or counted twice.

\begin{lemma}[the nine shifts]\label{lem:shifts}
The nine shifts $aV+bW$ $(a,b\in\{0,1,2\})$ are pairwise distinct. In increasing order they are
\begin{equation}\label{eq:shifts}
\begin{gathered}
e_0=0,\quad e_1=V,\quad e_2=2V,\quad e_3=2V+U,\quad e_4=3V+U,\\
e_5=4V+U,\quad e_6=4V+2U,\quad e_7=5V+2U,\quad e_8=6V+2U,
\end{gathered}
\end{equation}
with consecutive gaps $V,V,U,V,V,U,V,V$.
\end{lemma}

\begin{proof}
$aV+bW=(a+2b)V+bU$. For fixed $b$ the three values increase in $a$ by steps of $V$; the blocks
$b=0,1,2$ start at $0$, $W=2V+U$ and $2W=4V+2U$, and each block starts at distance $U>0$ above the
last element of the previous block.
\end{proof}

\begin{lemma}[window lemma]\label{lem:window}
For every $n\in\Z$, at most two shifts $e_i$ satisfy $0\le n-e_i\le D$. If two do, they are
consecutive in the order~\eqref{eq:shifts}.
\end{lemma}

\begin{proof}
The admissible shifts lie in the closed interval $[n-D,n]$, of length $D$. If three shifts
$e_i<e_j<e_k$ lay in it, so would the three consecutive shifts $e_i,e_{i+1},e_{i+2}$, whose span is
the sum of two consecutive gaps, namely $2V$ or $V+U$. Since $V>U$, both spans are at least
$V+U=D+3>D$, a contradiction. The same estimate excludes two non-consecutive shifts, whose span is
at least that of three consecutive ones.
\end{proof}

The margin is exactly $3$, so the conclusion survives the fact that the interval is closed. The
argument uses only the degree $D$ of $G_\ell$ and the gap pattern, not the shape of the coefficients
of $G_\ell$.

\begin{proposition}\label{prop:upper}
$B_{\ell+2}\le2B_\ell$ for every $\ell\ge1$.
\end{proposition}

\begin{proof}
Fix $n$ and apply~\eqref{eq:twostep}. A term $\rho_\ell(n-e_i)$ is non-zero only if
$0\le n-e_i\le D_\ell$, so by Lemma~\ref{lem:window} at most two terms are non-zero. Each is at most
$B_\ell$. Hence $\rho_{\ell+2}(n)\le2B_\ell$.
\end{proof}

\begin{thmPETone}\phantomsection\label{thm:pet1}
For every $\ell\ge1$,
\[
B_\ell=2^{\floor{(\ell-1)/2}}.
\]
Moreover, for every $\ell\ge1$ there is an integer $x_\ell$ with
\begin{equation}\label{eq:pair}
\rho_\ell(x_\ell)=\rho_\ell\bigl(x_\ell-P_{\ell+1}\bigr)=B_\ell .
\end{equation}
\end{thmPETone}

\begin{proof}
We prove~\eqref{eq:pair} and $B_{\ell+2}=2B_\ell$ simultaneously, by induction from $\ell$ to
$\ell+2$, separately on odd and on even $\ell$; we keep the notation $U,V,W,D$.

\emph{Bases.} For $\ell=1$, $G_1=1+x^2+x^4$, so $B_1=1$, and $x_1=2=P_2$ gives
$\rho_1(2)=\rho_1(0)=1$. For $\ell=2$ the nine sums $2a+5b$ $(a,b\in\{0,1,2\})$ are
$0,2,4,5,7,9,10,12,14$, pairwise distinct, so $B_2=1$, and $x_2=5=P_3$ gives
$\rho_2(5)=\rho_2(0)=1$.

\emph{Step.} Assume~\eqref{eq:pair} for $\ell$. By Lemma~\ref{lem:shifts}, $(e_2,e_3)=(2V,2V+U)$ and
$(e_5,e_6)=(4V+U,4V+2U)$ are pairs of consecutive shifts at distance $U$. For
$e\in\{2V,\,4V+U\}$ put $n=e+x_\ell$. Then~\eqref{eq:twostep} contains the two terms
\[
\rho_\ell(n-e)=\rho_\ell(x_\ell)=B_\ell,\qquad
\rho_\ell(n-e-U)=\rho_\ell(x_\ell-P_{\ell+1})=B_\ell,
\]
and every other term is non-negative. Hence $\rho_{\ell+2}(n)\ge2B_\ell$, and with
Proposition~\ref{prop:upper}, $B_{\ell+2}=2B_\ell$. Both targets $2V+x_\ell$ and $4V+U+x_\ell$ are
therefore maximal, and they differ by $2V+U=P_{\ell+3}=P_{(\ell+2)+1}$. So~\eqref{eq:pair} holds for
$\ell+2$ with $x_{\ell+2}=4V+U+x_\ell$.

\emph{Conclusion.} From $B_{\ell+2}=2B_\ell$ and $B_1=B_2=1$ we get $B_{2k-1}=B_{2k}=2^{k-1}$ for
$k\ge1$, which is $2^{\floor{(\ell-1)/2}}$.
\end{proof}

\begin{remark}[no circularity]\label{rem:noncirc}
The lower bound in the step uses only~\eqref{eq:pair} at length $\ell$ and the non-negativity of the
coefficients; it never uses $B_{\ell+2}$. The upper bound uses only the degree of $G_\ell$
(Proposition~\ref{prop:upper}). The pair at length $\ell+2$ is derived from the two bounds, not
assumed. In particular $B_{4m+1}=4^m$, the case used in Section~\ref{sec:stt2}.
\end{remark}

%======================================================================
\section{Odd-length extremal classification}\label{sec:pet2}
%======================================================================

\begin{thmPETtwo}\phantomsection\label{thm:pet2}
Let $\ell\ge1$ be odd, put $U=P_{\ell+1}$ and $u_\ell=(P_{\ell+2}-5)/2$. Then $u_\ell$ is an
integer, $D_\ell=2u_\ell+U+2$, and
\[
M_\ell=\{u_\ell,\ u_\ell+2,\ u_\ell+U,\ u_\ell+U+2\}
=\{u_\ell,\ u_\ell+2,\ D_\ell-u_\ell-2,\ D_\ell-u_\ell\}.
\]
For $\ell\ge3$ these are four distinct integers. For $\ell=1$ the set has three elements,
$M_1=\{0,2,4\}$.
\end{thmPETtwo}

\begin{proof}
Since $\ell+2$ is odd, $P_{\ell+2}$ is odd by Lemma~\ref{lem:sums}(b), so $u_\ell\in\Z$. By
Lemma~\ref{lem:degree}, $2u_\ell+U+2=P_{\ell+2}+P_{\ell+1}-3=D_\ell$; this also converts the two
descriptions of $M_\ell$ into each other, since $D_\ell-u_\ell-2=u_\ell+U$ and
$D_\ell-u_\ell=u_\ell+U+2$. For $\ell\ge3$, $U\ge P_4=12>2$, so the four numbers are distinct.
We argue by induction on odd $\ell$.

\emph{Base $\ell=1$.} Here $u_1=(P_3-5)/2=0$, $U=P_2=2$ and $D_1=4$. Directly from
$G_1=1+x^2+x^4$, $M_1=\{0,2,4\}$, which is the stated set (the elements $u_1+2$ and $u_1+U$
coincide).

\emph{Step.} Assume the claim at an odd $\ell$, write $B=B_\ell$, $u=u_\ell$, and use $U,V,W,D$ of
Section~\ref{sec:pet1}. By \PETone, $B_{\ell+2}=2B$. Recall that
\begin{equation}\label{eq:VbigU}
V=2U+P_\ell>U+2.
\end{equation}

\emph{Necessity.} Suppose $\rho_{\ell+2}(n)=2B$. By Lemma~\ref{lem:window} at most two terms
of~\eqref{eq:twostep} are non-zero, and each is at most $B$; since they sum to $2B$, there are
exactly two of them, they belong to consecutive shifts $e<e'$, and each equals $B$. Then
$n-e,\ n-e'\in M_\ell$, so $e'-e$ is a positive difference of two elements of $M_\ell$. The positive
differences of $\{u,u+2,u+U,u+U+2\}$ lie in
\begin{equation}\label{eq:diffs}
\{2,\ U-2,\ U,\ U+2\}.
\end{equation}
On the other hand $e'-e\in\{V,U\}$ by Lemma~\ref{lem:shifts}, and $V$ does not belong
to~\eqref{eq:diffs} by~\eqref{eq:VbigU}. Hence $e'-e=U$, and by Lemma~\ref{lem:shifts} the only
consecutive pairs at distance $U$ are $(2V,2V+U)$ and $(4V+U,4V+2U)$; thus $e\in\{2V,4V+U\}$.

Put $y=n-e'$. Then $n-e=y+U$, and $y,\ y+U\in M_\ell$. Every element of $M_\ell$ is at least $u$,
so $y+U\ge u+U$. The elements of $M_\ell$ that are at least $u+U$ are exactly $u+U$ and $u+U+2$:
this is clear for $U>2$, because $u<u+2<u+U<u+U+2$, and for $\ell=1$ (where $U=2$) the elements
$u+2$ and $u+U$ coincide and the list is still $\{u+U,u+U+2\}$. Hence $y\in\{u,u+2\}$, and
\[
n=e'+y\in\{\,2V+U+u,\ 2V+U+u+2,\ 4V+2U+u,\ 4V+2U+u+2\,\}.
\]

\emph{Sufficiency.} For each of these four $n$, the consecutive pair of shifts $(e,e+U)$ with
$e\in\{2V,4V+U\}$ contributes the two terms $B+B$. Hence $\rho_{\ell+2}(n)\ge2B=B_{\ell+2}$, and by
Proposition~\ref{prop:upper} each such $n$ is a maximizer.

\emph{Shape.} Put $u'=2V+U+u$ and $U'=P_{\ell+3}=2V+U$. The four points are
$u',\ u'+2,\ u'+U',\ u'+U'+2$. Since $P_{\ell+4}=2P_{\ell+3}+P_{\ell+2}=5V+2U$,
\[
u'=\frac{4V+2U+(V-5)}{2}=\frac{P_{\ell+4}-5}{2}=u_{\ell+2},
\]
which is the claim at length $\ell+2$.
\end{proof}

\begin{remark}[specialization to $\ell=4m+1$]\label{rem:special}
Take $\ell=4m+1$ and $s=2m+2$. Then $U=P_{4m+2}=P_{2s-2}$ (not $P_s$, unless $m=0$), and by
Lemma~\ref{lem:Cm} below $u_\ell=(P_{4m+3}-5)/2=z_m-2$, where $z_m=C_mP_s$. Using
$u_\ell+U=D_\ell-u_\ell-2$ and $H_m=D_{4m+1}$,
\[
\argmax_nR_m(n)=\{\,z_m-2,\ z_m,\ H_m-z_m,\ H_m-z_m+2\,\}.
\]
For $m=0$ this is $\{0,2,4\}$; for $m\ge1$ it consists of four distinct integers.
\end{remark}

%======================================================================
\section{Even-length recurrence and counting}\label{sec:pet3}
%======================================================================

For a length $L\ge1$ and an integer $g$ put
\[
T_L(g)=\{x\in M_L:\ x-g\in M_L\}.
\]

\begin{lemma}[extremal-set recurrence]\label{lem:setrec}
For every $\ell\ge1$, with the shifts $e_0<\dots<e_8$ of~\eqref{eq:shifts},
\[
M_{\ell+2}=\bigcup_{i=0}^{7}\bigl((M_\ell+e_i)\cap(M_\ell+e_{i+1})\bigr)
=\bigcup_{i=0}^{7}\bigl(e_i+T_\ell(e_{i+1}-e_i)\bigr),
\]
and the eight pieces are pairwise disjoint.
\end{lemma}

\begin{proof}
By \PETone, $B_{\ell+2}=2B_\ell$. By Lemma~\ref{lem:window}, $\rho_{\ell+2}(n)=2B_\ell$ requires two
non-zero summands in~\eqref{eq:twostep}, both equal to $B_\ell$ and belonging to consecutive shifts
$e_i,e_{i+1}$, i.e.\ $n\in(M_\ell+e_i)\cap(M_\ell+e_{i+1})$. Conversely, a point of such an
intersection has two summands equal to $B_\ell$ and, by Lemma~\ref{lem:window}, no third non-zero
summand, so its coefficient is $2B_\ell$. A point common to two different pieces would have three
non-zero summands, which Lemma~\ref{lem:window} excludes. Finally
$(M_\ell+e_i)\cap(M_\ell+e_{i+1})=e_i+T_\ell(e_{i+1}-e_i)$ by definition of $T_\ell$.
\end{proof}

From now on let $\ell=2k$ with $k\ge1$, and put
\[
a=P_{2k+2},\qquad c=P_{2k+1},\qquad e=P_{2k+3}=2a+c,\qquad v=\frac{c-5}{2}.
\]
In the notation of Section~\ref{sec:pet1} (with $\ell=2k$), $V=a$, $U=c$, $W=e$; $v$ is an integer
because $c$ is odd. (Here $e$ denotes the number $P_{2k+3}$, not one of the shifts $e_i$
of~\eqref{eq:shifts}.) Let
\[
X_k=T_{2k}(a),\qquad Y_k=T_{2k}(c),\qquad q_{2k}=|Y_k|.
\]

\begin{lemma}[even blocks]\label{lem:blocks}
For every $k\ge1$,
\[
M_{2k+2}=(Q_X+X_k)\sqcup(Q_Y+Y_k),\qquad Q_X=\{0,a,e,a+e,2e,a+2e\},\qquad Q_Y=\{2a,\,2a+e\},
\]
where the six translates $h+X_k$ $(h\in Q_X)$ and the two translates $r+Y_k$ $(r\in Q_Y)$ are
pairwise disjoint.
\end{lemma}

\begin{proof}
In the order~\eqref{eq:shifts} the shifts are $0,a,2a,e,a+e,2a+e,2e,a+2e,2a+2e$, with gaps
$a,a,c,a,a,c,a,a$. The two gaps equal to $c$ start at $e_2=2a$ and $e_5=2a+e$; the six gaps equal to
$a$ start at $0,a,e,a+e,2e,a+2e$. Apply Lemma~\ref{lem:setrec}.
\end{proof}

\begin{lemma}[$X_k$ and the extreme maximizers]\label{lem:X}
For every $k\ge1$: $\min M_{2k}=v$, $\max M_{2k}=a+v+2$, $v+2\in M_{2k}$, $v+1\notin M_{2k}$,
and
\[
X_k=\{a+v,\ a+v+2\};\qquad\text{in particular }|X_k|=2.
\]
\end{lemma}

\begin{proof}
Let $L=2k$, so $L-1$ is odd. By \PETtwo{} at length $L-1$ (where
$u_{L-1}=(P_{L+1}-5)/2=v$ and $U=P_L$), $M_{L-1}=\{v,v+2,v+P_L,v+P_L+2\}$; hence
$\min M_{L-1}=v$, $v+2\in M_{L-1}$ and $v+1\notin M_{L-1}$. Since $G_L=G_{L-1}(1+x^c+x^{2c})$, only
the constant term of the last factor contributes to coefficients below $c$, so
$\rho_L(n)=\rho_{L-1}(n)$ for $n<c$; and $B_L=B_{L-1}$ by \PETone, because
$\floor{(L-1)/2}=\floor{(L-2)/2}$. Hence $M_L\cap[0,c)=M_{L-1}\cap[0,c)$, and since $v+2=(c-1)/2<c$
we obtain $\min M_L=v$, $v+2\in M_L$ and $v+1\notin M_L$. By Lemma~\ref{lem:degree},
$D_L=a+c-3=a+2v+2$, so by symmetry $\max M_L=D_L-v=a+v+2$ and $a+v=D_L-(v+2)\in M_L$.
If $x\in X_k$, then $x-a\ge\min M_L=v$ and $x\le\max M_L=a+v+2$, so $x-a\in\{v,v+1,v+2\}$; as
$v+1\notin M_L$, only $x-a\in\{v,v+2\}$ remains, and both $x=a+v$ and $x=a+v+2$ belong to $X_k$
because $a+v,\ a+v+2,\ v,\ v+2\in M_L$.
\end{proof}

\begin{lemma}[the $Y$-recurrence]\label{lem:Y}
$Y_1=\{5,7,9,10,12,14\}$, so $q_2=6$. For every $k\ge1$, every $y\in Y_k$ satisfies
$c+v\le y\le a+v+2$, and
\[
Y_{k+1}=\bigl(\{e,\,a+e,\,2e,\,a+2e\}+X_k\bigr)\ \sqcup\ \bigl(2a+e+Y_k\bigr).
\]
Hence $q_{2k+2}=q_{2k}+8$ and $q_{2k}=8k-2$ for every $k\ge1$.
\end{lemma}

\begin{proof}
For $k=1$, $M_2=\{0,2,4,5,7,9,10,12,14\}$ (proof of \PETone) and $c=5$; the elements $y$ with
$y-5\in M_2$ are $5,7,9,10,12,14$.

\emph{Bounds.} If $y\in Y_k$ then $y-c\in M_{2k}$, so $y\ge c+\min M_{2k}=c+v$, and
$y\le\max M_{2k}=a+v+2$ by Lemma~\ref{lem:X}. In particular every difference of two elements of
$Y_k$ is at most $a+2-c<e$. Since $X_k=\{a+v,a+v+2\}$ and $y\ge v$, every $x\in X_k$ and
$y\in Y_k$ satisfy $-2\le x-y\le a+2$.

\emph{Classification.} By definition $Y_{k+1}=T_{2k+2}(e)=\{z\in M_{2k+2}:z-e\in M_{2k+2}\}$. By
Lemma~\ref{lem:blocks} every element of $M_{2k+2}$ lies in exactly one block $h+X_k$ $(h\in Q_X)$
or $r+Y_k$ $(r\in Q_Y)$. We classify the pairs $(z,z-e)$ according to the blocks of $z$ and of
$z-e$.

\emph{$X$--$X$.} Let $z=h''+x''$ and $z-e=h'+x'$ with $h',h''\in Q_X$ and $x',x''\in X_k$. Then
$x''-x'\in\{-2,0,2\}$, so $h''-h'\in\{e-2,e,e+2\}$. The positive differences of $Q_X$ are
$a,\ e-a=a+c,\ e,\ a+e,\ 2e-a=3a+2c,\ 2e,\ a+2e$, and all except $e$ differ from $e$ by at least
$a\ge12>2$. Hence $h''-h'=e$ and $x''=x'$. The pairs $(h',h'')$ with $h''-h'=e$ are $(0,e)$,
$(a,a+e)$, $(e,2e)$ and $(a+e,a+2e)$, so this case contributes exactly
$\{e,a+e,2e,a+2e\}+X_k$, eight points.

\emph{$Y$--$Y$.} Let $z=r''+y''$ and $z-e=r'+y'$ with $r',r''\in Q_Y=\{2a,2a+e\}$, so that
$(r''-r')+(y''-y')=e$ and $r''-r'\in\{-e,0,e\}$. The cases $r''=r'$ and $r''=r'-e$ would require
$y''-y'\in\{e,2e\}$, which is impossible because $|y''-y'|<e$. Hence $r'=2a$, $r''=2a+e$ and
$y''=y'$. This contributes exactly $2a+e+Y_k$.

\emph{$X$--$Y$.} Let $z=\xi+x$ with $\xi\in Q_X$, $x\in X_k$, and $z-e=\eta+y$ with $\eta\in Q_Y$,
$y\in Y_k$. Then $x-y=e-(\xi-\eta)$. The eight values of $\xi-\eta$ are
$-4a-c,\,-3a-c,\,-2a,\,-a,\,c,\,a+c,\,2a+2c,\,3a+2c$, so the required values of $x-y$ are
$6a+2c,\,5a+2c,\,4a+c,\,3a+c,\,2a,\,a,\,-c,\,-a-c$. Since $-2\le x-y\le a+2<2a$ and
$-c\le-5<-2$, only $x-y=a$ is possible. It forces $y=x-a\in\{v,v+2\}$, whereas $y\ge c+v>v+2$.
No such pair exists.

\emph{$Y$--$X$.} Let $z=\eta+y$ and $z-e=\xi+x$. Then $y-x=e-(\eta-\xi)$ takes the eight values
$-2a,\,-a,\,c,\,a+c,\,2a+2c,\,3a+2c,\,4a+3c,\,5a+3c$. Since $-a-2\le y-x\le2$, $-2a<-a-2$ and
$c>2$, only $y-x=-a$ is possible. It forces $y=x-a\in\{v,v+2\}$ again. No such pair exists.

The two surviving parts are disjoint, because the first lies in $Q_X+X_k$ and the second in
$Q_Y+Y_k$. Counting gives $q_{2k+2}=4\cdot|X_k|+q_{2k}=q_{2k}+8$, and with $q_2=6$, $q_{2k}=8k-2$.
\end{proof}

\begin{thmPETthree}\phantomsection\label{thm:pet3}
Let $E_k=|M_{2k}|$. Then $E_1=9$, and $E_k=16k-8$ for every $k\ge2$. The formula $16k-8$ does not
hold at $k=1$.
\end{thmPETthree}

\begin{proof}
By the proof of \PETone, $B_2=1$ and $M_2=\{0,2,4,5,7,9,10,12,14\}$, so $E_1=9$. For $k\ge1$,
Lemma~\ref{lem:blocks} decomposes $M_{2k+2}$ into six pairwise disjoint translates of $X_k$ and two
pairwise disjoint translates of $Y_k$, so by Lemmas~\ref{lem:X} and~\ref{lem:Y}
\[
E_{k+1}=6\cdot2+2(8k-2)=16(k+1)-8 .
\]
At $k=1$ the formula would give $8\ne9$.
\end{proof}

%======================================================================
\section{Explicit parametrization}\label{sec:pet4}
%======================================================================

For $k\ge1$ and $j\ge3$ put
\[
S_j=P_3+P_4+\dots+P_j,\qquad K=\{0,2,4,5,7,9\},\qquad z_1=\frac{P_{2k+2}}{2}-6,\qquad
z_2=z_1+P_{2k+1},
\]
and define the six families
\[
\begin{aligned}
\CC_1&=z_1+K, &\qquad \CC_2&=z_2+K,\\
\OL_j&=\{z_1-S_j,\ z_1-S_j+2\}, && 3\le j\le 2k,\\
\IR_j&=\{z_1+7+S_j,\ z_1+9+S_j\}, && 3\le j\le 2k-1,\\
\IL_j&=\{z_2-S_j,\ z_2-S_j+2\}, && 3\le j\le 2k-1,\\
\OR_j&=\{z_2+7+S_j,\ z_2+9+S_j\}, && 3\le j\le 2k.
\end{aligned}
\]
We call $\CC_1,\CC_2$ the \emph{cores}, $\OL_j,\OR_j$ the \emph{outer chains} and $\IR_j,\IL_j$ the
\emph{inner chains}, and we write
\[
\Fam_k=\CC_1\cup\CC_2\cup\bigcup_{j=3}^{2k}(\OL_j\cup\OR_j)\cup\bigcup_{j=3}^{2k-1}(\IR_j\cup\IL_j).
\]
(Sans-serif letters denote these families; $C_m$ in italics is the silver-ratio target of
Section~\ref{sec:silvermodel}.) For $k=1$ all four chains are empty.

\begin{thmPETfour}\phantomsection\label{thm:pet4}
For every $k\ge1$, $M_{2k}=\Fam_k$. If $k\ge2$, the $16k-8$ numbers listed in the six families are
pairwise distinct, and blockwise
\[
\OL_{2k}<\dots<\OL_3<\CC_1<\IR_3<\dots<\IR_{2k-1}<\IL_{2k-1}<\dots<\IL_3<\CC_2<\OR_3<\dots<\OR_{2k},
\]
so $|M_{2k}|=16k-8$. If $k=1$, then $z_1=0$, $z_2=5$, the chains are empty,
$\CC_1\cap\CC_2=\{5,7,9\}$ and $M_2=\CC_1\cup\CC_2$ has $9$ elements.
\end{thmPETfour}

The proof occupies the rest of this section. We keep the notation $a,c,e,v,X_k,Y_k$ of
Section~\ref{sec:pet3} at level $k$, and we write $z_1',z_2'$ and $\CC_1',\OL_j',\dots$ for the
corresponding objects at level $k+1$; thus $z_1'=P_{2k+4}/2-6$ and $z_2'=z_1'+P_{2k+3}=z_1'+e$.

\begin{lemma}[arithmetic of the parameters]\label{lem:arith}
\begin{enumerate}[label=\textup{(\roman*)}]
\item $S_j=(P_j+P_{j+1}-7)/2$ is an integer, and $S_{j+1}-S_j=P_{j+1}\ge12$ for $j\ge3$.
\item $z_1=\sum_{j=2}^{k}P_{2j+1}$ (an empty sum for $k=1$); in particular $z_1,z_2\in\Z$.
\item $D_{2k}=z_1+z_2+9$.
\item For $k\ge2$, $S_{2k-1}+S_{2k}+7=c$.
\end{enumerate}
\end{lemma}

\begin{proof}
(i) By Lemma~\ref{lem:sums}(a), $S_j=(P_{j+1}+P_j-1)/2-P_1-P_2=(P_j+P_{j+1}-7)/2$; the numerator is
even because $P_j$ and $P_{j+1}$ have different parities (Lemma~\ref{lem:sums}(b)).
(ii) By Lemma~\ref{lem:sums}(c), $\sum_{i=0}^{k}P_{2i+1}=P_{2k+2}/2$; subtract $P_1+P_3=6$.
(iii) $z_1+z_2+9=2z_1+c+9=a+c-3=D_{2k}$ by Lemma~\ref{lem:degree}.
(iv) $S_{2k-1}+S_{2k}+7=(P_{2k-1}+2P_{2k}+P_{2k+1})/2=(2P_{2k+1})/2=c$.
\end{proof}

\begin{lemma}[nine identities]\label{lem:nine}
For every $k\ge1$:
\begin{align*}
&(1)\quad z_1'=z_1+e, &&(2)\quad z_2'=z_1+2e=z_2+2a+e,\\
&(3)\quad S_{2k-1}+S_{2k}+7=c\quad(k\ge2), &&(4)\quad z_1'+7+S_{2k+1}=e+a+v,\\
&(5)\quad z_2'-S_{2k+1}=e+2a+v, &&(6)\quad z_2'+7+S_{2k+1}=2e+a+v,\\
&(7)\quad z_2'+7+S_{2k+2}=2e+2a+v, &&(8)\quad z_1'-S_{2k+1}=2a+v,\\
&(9)\quad z_1'-S_{2k+2}=a+v.
\end{align*}
\end{lemma}

\begin{proof}
We use $P_{2k+4}=2e+a$, $e=2a+c$, $S_{2k+1}=(a+c-7)/2$, $S_{2k+2}=(a+e-7)/2$ (Lemma~\ref{lem:arith}(i))
and $(e-5)/2=a+v$.
(1) $z_1'=(2e+a)/2-6=e+z_1$.
(2) $z_2'=z_1'+e=z_1+2e$, and $z_2+2a+e=z_1+c+2a+e=z_1+2e$.
(3) is Lemma~\ref{lem:arith}(iv).
(4) $z_1'+7+S_{2k+1}=e+\tfrac a2+1+\tfrac{a+c-7}{2}=e+a+\tfrac{c-5}{2}=e+a+v$.
(5) $z_2'-S_{2k+1}=(z_1'+7+S_{2k+1})+e-7-2S_{2k+1}=2e+a+v-(a+c)=e+2a+v$, using $e-c=2a$.
(6) $z_2'+7+S_{2k+1}=(z_1'+7+S_{2k+1})+e=2e+a+v$ by (4).
(7) $z_2'+7+S_{2k+2}=2e+\tfrac a2+1+\tfrac{a+e-7}{2}=2e+a+\tfrac{e-5}{2}=2e+2a+v$.
(8) $z_1'-S_{2k+1}=(z_2'-S_{2k+1})-e=2a+v$ by (5).
(9) $z_1'-S_{2k+2}=e+\tfrac a2-6-\tfrac{a+e-7}{2}=\tfrac{e-5}{2}=a+v$.
\end{proof}

For $k\ge1$ let
\[
\Gamma_k=\CC_2\cup\bigcup_{j=3}^{2k-1}\IL_j\cup\bigcup_{j=3}^{2k}\OR_j\cup\IR_{2k-1},
\]
with the families taken at level $k$, where the last term is present only for $k\ge2$.

\begin{proposition}[closed form of $Y_k$]\label{prop:Yclosed}
$Y_k=\Gamma_k$ for every $k\ge1$.
\end{proposition}

\begin{proof}
For $k=1$, $z_2=5$ and all chains are empty, so $\Gamma_1=5+K=Y_1$ by Lemma~\ref{lem:Y}. Assume
$Y_k=\Gamma_k$ for some $k\ge1$. By Lemma~\ref{lem:Y} it suffices to show
\[
\Gamma_{k+1}=\bigl(\{e,a+e,2e,a+2e\}+X_k\bigr)\cup\bigl(2a+e+\Gamma_k\bigr).
\]
\emph{The translate.} By identity~(2), $z_2+2a+e=z_2'$, so translation by $2a+e$ maps $\CC_2$ onto
$\CC_2'$, $\IL_j$ onto $\IL_j'$ $(3\le j\le2k-1)$ and $\OR_j$ onto $\OR_j'$ $(3\le j\le2k)$. For
$k\ge2$ it also maps $\IR_{2k-1}$ onto $\IL_{2k}'$: indeed, by identity~(3),
$z_1+7+S_{2k-1}+2a+e=z_1+2a+e+c-S_{2k}=z_1+2e-S_{2k}=z_2'-S_{2k}$ by identity~(2). Hence
\[
2a+e+\Gamma_k=\CC_2'\cup\bigcup_{j=3}^{2k}\bigl(\IL_j'\cup\OR_j'\bigr).
\]
\emph{The remaining families.} The families of $\Gamma_{k+1}$ not yet obtained are $\IL_{2k+1}'$,
$\OR_{2k+1}'$, $\OR_{2k+2}'$ and $\IR_{2k+1}'$. Each has the form $\{p,p+2\}$, and by
identities~(5), (6), (7) and~(4) their smallest elements $p$ are $e+2a+v$, $2e+a+v$, $2e+2a+v$
and $e+a+v$, that is, the four numbers $\{e,a+e,2e,a+2e\}+(a+v)$. Since $X_k=\{a+v,a+v+2\}$
(Lemma~\ref{lem:X}), these four families together are exactly $\{e,a+e,2e,a+2e\}+X_k$. This proves
the displayed identity, hence $Y_{k+1}=\Gamma_{k+1}$.
\end{proof}

\begin{proof}[Proof of $M_{2k}=\Fam_k$]
For $k=1$, $\Fam_1=K\cup(5+K)=\{0,2,4,5,7,9,10,12,14\}=M_2$. Let $k\ge1$; we show
$\Fam_{k+1}=M_{2k+2}$. By Lemma~\ref{lem:blocks} and Proposition~\ref{prop:Yclosed},
\[
Q_Y+Y_k=(2a+e+\Gamma_k)\cup\bigl((2a+e+\Gamma_k)-e\bigr).
\]
The first set was computed in the proof of Proposition~\ref{prop:Yclosed}. Subtracting $e$ replaces
$z_2'$ by $z_2'-e=z_1'$, so it maps $\CC_2'\mapsto\CC_1'$, $\IL_j'\mapsto\OL_j'$ and
$\OR_j'\mapsto\IR_j'$. Therefore
\[
Q_Y+Y_k=\CC_1'\cup\CC_2'\cup\bigcup_{j=3}^{2k}\bigl(\OL_j'\cup\IR_j'\cup\IL_j'\cup\OR_j'\bigr).
\]
The families of $\Fam_{k+1}$ not yet obtained are $\OL_{2k+1}'$, $\OL_{2k+2}'$, $\IR_{2k+1}'$,
$\IL_{2k+1}'$, $\OR_{2k+1}'$ and $\OR_{2k+2}'$. Each has the form $\{p,p+2\}$, and by
identities~(8), (9), (4), (5), (6) and~(7) their smallest elements are
$2a+v$, $a+v$, $e+a+v$, $e+2a+v$, $2e+a+v$ and $2e+2a+v$, that is, exactly $Q_X+(a+v)$. With
$X_k=\{a+v,a+v+2\}$ these six families are exactly $Q_X+X_k$. Hence
$\Fam_{k+1}=(Q_X+X_k)\cup(Q_Y+Y_k)=M_{2k+2}$ by Lemma~\ref{lem:blocks}.
\end{proof}

\begin{remark}[no circularity]\label{rem:noncirc4}
Proposition~\ref{prop:Yclosed} uses only Lemmas~\ref{lem:X}, \ref{lem:Y} and~\ref{lem:nine}, never
the identity $M_{2k}=\Fam_k$; the proof of $M_{2k}=\Fam_k$ then uses
Proposition~\ref{prop:Yclosed}. No comparison of cardinalities is used: $M_{2k}=\Fam_k$ is proved as
an identity of sets.
\end{remark}

\begin{proposition}[ordering and count]\label{prop:order}
Let $k\ge2$. The six families are pairwise disjoint, the blocks are ordered as displayed in
\PETfour, and $|\Fam_k|=16k-8$.
\end{proposition}

\begin{proof}
Within a chain, consecutive smallest elements differ by $S_{j+1}-S_j=P_{j+1}\ge12>2$
(Lemma~\ref{lem:arith}(i)), so the two-element blocks of a chain are disjoint and ordered by $j$
(decreasing in $\OL$ and $\IL$, increasing in $\IR$ and $\OR$). Between the families:
$\max\OL=z_1-S_3+2=z_1-3<z_1=\min\CC_1$; $\max\CC_1=z_1+9<z_1+12=z_1+7+S_3=\min\IR$;
$\max\IR=z_1+9+S_{2k-1}<z_2-S_{2k-1}=\min\IL$, which is equivalent to $2S_{2k-1}+9<c$, i.e.\ by
Lemma~\ref{lem:arith}(iv) to $S_{2k-1}+2<S_{2k}$, i.e.\ to $2<P_{2k}$, true for $k\ge2$;
$\max\IL=z_2-S_3+2=z_2-3<z_2=\min\CC_2$; and $\max\CC_2=z_2+9<z_2+12=\min\OR$. Counting,
$|\OL|=|\OR|=2(2k-2)$, $|\IR|=|\IL|=2(2k-3)$ and $|\CC_1|=|\CC_2|=6$, so
$|\Fam_k|=2(4k-4)+2(4k-6)+12=16k-8$.
\end{proof}

\begin{proof}[Completion of the proof of \PETfour]
The identity $M_{2k}=\Fam_k$ was proved above, and Proposition~\ref{prop:order} gives the ordering,
disjointness and count for $k\ge2$. For $k=1$, $z_1=P_4/2-6=0$ and $z_2=5$; the index ranges of all
four chains are empty, $\CC_1=\{0,2,4,5,7,9\}$, $\CC_2=\{5,7,9,10,12,14\}$, $\CC_1\cap\CC_2=\{5,7,9\}$,
and the twelve listed numbers collapse to the nine elements of $M_2$.
\end{proof}

\begin{remark}[extremes and symmetry]\label{rem:symmetry4}
For $k\ge2$, $\min M_{2k}=\min\OL_{2k}=z_1-S_{2k}=(a-P_{2k}-c-5)/2=v$, because $a=2c+P_{2k}$, and
$\max M_{2k}=\max\OR_{2k}=z_2+9+S_{2k}=D_{2k}-v$, in agreement with Lemma~\ref{lem:X}. By
Lemma~\ref{lem:arith}(iii) the reflection $n\mapsto D_{2k}-n=z_1+z_2+9-n$ of Lemma~\ref{lem:degree}
maps $\CC_1$ onto $\CC_2$ (because $9-K=K$), $\OL_j$ onto $\OR_j$, and $\IR_j$ onto $\IL_j$:
\[
D_{2k}-(z_1+7+S_j)=z_2-S_j+2,\qquad D_{2k}-(z_1+9+S_j)=z_2-S_j .
\]
The index ranges $3\le j\le2k$ of $\OL,\OR$ and $3\le j\le2k-1$ of $\IR,\IL$ match, as the symmetry
requires.
\end{remark}

\begin{remark}[scope]\label{rem:scope4}
Together with Proposition~\ref{prop:order}, \PETfour{} gives a second derivation of the count in
\PETthree; it uses the set identities of Lemmas~\ref{lem:blocks}--\ref{lem:Y} but not the cardinality
recursion $q_{2k+2}=q_{2k}+8$. \PETfour{} is a description of $M_{2k}$; we make no claim that it is
the simplest possible one, and it says nothing about coefficients close to the maximum.
\end{remark}

%======================================================================
\section{Silver-ratio transfer}\label{sec:stt1}
%======================================================================

\begin{thmSTTone}\phantomsection\label{thm:stt1}
For every $m\ge0$ and every $t\in\Z$, with $s=2m+2$,
\[
\nu_m(t)=R_m(tP_s)=\rho_{4m+1}(tP_{2m+2}).
\]
The identity is induced by the bijection $(d_k)_{k=-2m}^{2m}\mapsto(d_{j-s})_{j=2}^{4m+2}$ between the
words of the two models; it is not an estimate.
\end{thmSTTone}

\begin{proof}
Let $d=(d_k)_{k=-2m}^{2m}\in\{0,1,2\}^{4m+1}$ and put $X=\sum_kd_k\delta^k-t$. By
Lemma~\ref{lem:powers},
\[
\delta^sX=A\,\delta+B',\qquad A=\sum_kd_kP_{k+s}-tP_s,\qquad B'=\sum_kd_kP_{k+s-1}-tP_{s-1}.
\]
Since $\{1,\delta\}$ is a basis of $\Q(\sqrt2)$ over $\Q$, $X=0$ if and only if $A=B'=0$.

\emph{Step 1 (shift).} As $k$ runs over $\{-2m,\dots,2m\}$, the index $k+s$ runs over
$\{2,\dots,4m+2\}$, which is the weight set of $G_{4m+1}$ in~\eqref{eq:model}.

\emph{Step 2 (the second coordinate is redundant).} For even $s$ and every $k\in\Z$,
\begin{equation}\label{eq:redund}
P_{k+s-1}P_s-P_{k+s}P_{s-1}=-P_k.
\end{equation}
Indeed, as functions of $k$ both sides satisfy the Pell recurrence (the left side is a fixed linear
combination of $P_{k+s-1}$ and $P_{k+s}$). They agree at $k=0$, where both vanish, and at $k=1$,
where the left side is $P_s^2-P_{s-1}P_{s+1}=(-1)^{s+1}=-1=-P_1$ by Lemma~\ref{lem:cassini}, since
$s$ is even. Multiplying $B'$ by $P_s$ and $A$ by $P_{s-1}$ and subtracting gives
\[
P_sB'-P_{s-1}A=-\sum_kd_kP_k=:-T .
\]

\emph{Step 3 (a strict bound on $T$).} By $P_{-n}=(-1)^{n+1}P_n$, the coefficient of $d_k$ in $T$ is
$P_k\ge0$ for $k\ge0$ and for negative odd $k$, and it is $-P_{|k|}$ for negative even $k$. Hence
\[
-2\sum_{i=1}^{m}P_{2i}\ \le\ T\ \le\ 2\sum_{k=1}^{2m}P_k+2\sum_{i=1}^{m}P_{2i-1}.
\]
By Lemma~\ref{lem:sums}(a) and~(c), $2\sum_{i\le m}P_{2i}=P_{2m+1}-1$,
$2\sum_{i\le m}P_{2i-1}=P_{2m}$ and $2\sum_{k\le2m}P_k=P_{2m+1}+P_{2m}-1$, so
\[
-(P_{2m+1}-1)\le T\le P_{2m+1}+2P_{2m}-1 .
\]
Since $P_s=2P_{2m+1}+P_{2m}$, we have $P_s-(P_{2m+1}+2P_{2m}-1)=P_{2m+1}-P_{2m}+1>0$ and
$-P_s<-(P_{2m+1}-1)$. Hence $|T|<P_s$.

\emph{Step 4 (conclusion).} If $A=0$, then $P_sB'=-T$, so $P_s\mid T$; as $|T|<P_s$, this forces
$T=0$ and then $B'=0$. Thus $X=0$ holds exactly when $A=0$, that is, when
$\sum_kd_kP_{k+s}=tP_s$. By Step~1 this is the Pell equation $\sum_{j=2}^{4m+2}d_{j-s}P_j=tP_s$ for
the word $(d_{j-s})_j$ of $G_{4m+1}$, and the shift of indices is a bijection between the two sets of
words.
\end{proof}

\begin{remark}[$m=0$]\label{rem:m0}
Here $s=2$ and the window is the single exponent $k=0$; the statement reads
$\nu_0(t)=\rho_1(2t)$, which equals $1$ for $t\in\{0,1,2\}$ and $0$ otherwise. Consistently,
Step~3 gives $0\le T\le0$.
\end{remark}

\begin{remark}[what \STTone{} is not]\label{rem:notstt1}
\STTone{} is an identity between two counting functions on the same window, written in two
coordinate systems. It is not an asymptotic statement, not a bound, and not a statement about other
windows, other digit alphabets or other quadratic units.
\end{remark}

%======================================================================
\section{Maximum multiplicity and uniqueness}\label{sec:stt2}
%======================================================================

\begin{lemma}\label{lem:Cm}
For every $m\ge0$, $C_m=P_{2m+1}+P_{2m}\in\Z$, and
\[
z_m:=C_mP_s=\frac{P_{4m+3}-1}{2}=\sum_{r=1}^{2m+1}P_{2r}.
\]
\end{lemma}

\begin{proof}
Summing the geometric series and using $\delta^2-1=2\delta$,
\[
C_m=\delta^{-2m}\,\frac{\delta^{4m+2}-1}{\delta^2-1}=\frac{\delta^{2m+2}-\delta^{-2m}}{2\delta}
=\frac{\delta^{2m+1}-\delta^{-(2m+1)}}{2}.
\]
Since $\lambda=-\delta^{-1}$ and $2m+1$ is odd, $\lambda^{2m+1}=-\delta^{-(2m+1)}$, so
$C_m=(\delta^{2m+1}+\lambda^{2m+1})/2=P_{2m+1}+P_{2m}$ by Lemma~\ref{lem:powers}. For even $s$ we have
$\lambda^s=\delta^{-s}$, so $P_s=(\delta^s-\delta^{-s})/(2\sqrt2)$; for odd $r$ we have
$\lambda^r=-\delta^{-r}$, so $P_r=(\delta^r+\delta^{-r})/(2\sqrt2)$. With $s=2m+2$,
\begin{align*}
C_mP_s&=\frac{(\delta^{2m+1}-\delta^{-(2m+1)})(\delta^{2m+2}-\delta^{-(2m+2)})}{4\sqrt2}\\
&=\frac{\delta^{4m+3}+\delta^{-(4m+3)}-(\delta+\delta^{-1})}{4\sqrt2}
=\frac{2\sqrt2\,P_{4m+3}-2\sqrt2\,P_1}{4\sqrt2}=\frac{P_{4m+3}-1}{2}.
\end{align*}
The last expression equals $\sum_{r=1}^{2m+1}P_{2r}$ by Lemma~\ref{lem:sums}(c).
\end{proof}

\begin{lemma}\label{lem:residue}
For every $m\ge0$, $H_m\equiv-4\pmod{P_s}$.
\end{lemma}

\begin{proof}
By Lemma~\ref{lem:degree}, $H_m=P_{4m+3}+P_{4m+2}-3=P_{2s-1}+P_{2s-2}-3$. The addition identity of
Lemma~\ref{lem:cassini} with $a=b=s$ gives $P_{2s}=P_s(P_{s+1}+P_{s-1})\equiv0\pmod{P_s}$, and the
recurrence $P_{2s}=2P_{2s-1}+P_{2s-2}$ then gives $P_{2s-2}\equiv-2P_{2s-1}$. With $a=s$, $b=s-1$
it gives $P_{2s-1}=P_s^2+P_{s-1}^2\equiv P_{s-1}^2$. Cassini's identity
$P_s^2-P_{s-1}P_{s+1}=(-1)^{s+1}=-1$, together with $P_{s+1}=2P_s+P_{s-1}\equiv P_{s-1}$, gives
$P_{s-1}^2\equiv1$. Hence $H_m\equiv1-2-3=-4$.
\end{proof}

\begin{thmSTTtwo}\phantomsection\label{thm:stt2}
For every $m\ge0$,
\[
N_m=\max_{t\in\Z}\nu_m(t)=4^m,
\qquad
\argmax_{t\in\Z}\nu_m(t)=
\begin{cases}
\{0,1,2\}, & m=0,\\
\{C_m\}, & m\ge1.
\end{cases}
\]
\end{thmSTTtwo}

\begin{proof}
By \STTone{} and \PETone, $\nu_m(t)=R_m(tP_s)\le B_{4m+1}=4^m$ for every $t\in\Z$.

\emph{Attainment for $m\ge1$.} Let $\ell=4m+1$. By Lemma~\ref{lem:Cm},
$u_\ell=(P_{4m+3}-5)/2=z_m-2$, so by \PETtwo{} (see Remark~\ref{rem:special}) the maximizers of
$R_m$ are $z_m-2$, $z_m$, $H_m-z_m$ and $H_m-z_m+2$, each with value $4^m$. In particular
$\nu_m(C_m)=R_m(C_mP_s)=R_m(z_m)=4^m$.

\emph{Uniqueness for $m\ge1$.} Let $t\in\Z$ with $\nu_m(t)=4^m$. By \STTone, $tP_s$ is one of the four
maximizers of $R_m$, and it is divisible by $P_s$. By Lemma~\ref{lem:residue} and $z_m\equiv0$, their
residues modulo $P_s$ are
\[
z_m-2\equiv-2,\qquad z_m\equiv0,\qquad H_m-z_m\equiv-4,\qquad H_m-z_m+2\equiv-2 .
\]
For $m\ge1$, $P_s\ge P_4=12>4$, so $-2$ and $-4$ are non-zero modulo $P_s$. Hence $tP_s=z_m$ and
$t=C_m$.

\emph{The case $m=0$.} Here $P_s=P_2=2$ and the maximizers of $R_0=\rho_1$ are $0,2,4$, all divisible
by $2$, so the residue argument gives no information: it needs $P_s>4$. By Remark~\ref{rem:m0},
$\nu_0(t)=1$ exactly for $t\in\{0,1,2\}$, so $N_0=1=4^0$ and $\argmax\nu_0=\{0,1,2\}$. The case $m=0$
is genuinely separate and is not a degenerate instance of the uniqueness statement.
\end{proof}

\begin{remark}[separation of the ingredients]\label{rem:sep}
Attainment and uniqueness are proved separately. Attainment uses \PETtwo{} at $\ell=4m+1$;
uniqueness uses the completeness part of \PETtwo{} together with Lemma~\ref{lem:residue}.
\end{remark}

%======================================================================
\section{General algebraic fibre bound}\label{sec:stt3}
%======================================================================

\begin{thmSTTthree}\phantomsection\label{thm:stt3}
Fix $m\ge0$ and the window $\{-2m,\dots,2m\}$. For every $\xi\in\Z[\sqrt2]$, the number of strings
$d\in\{0,1,2\}^{4m+1}$ with $\sum_kd_k\delta^k=\xi$ is at most $4^m$. The bound is attained.
\end{thmSTTthree}

\begin{proof}
Let $d$ be a string with $\sum_kd_k\delta^k=\xi$, and let $s=2m+2$. By Lemma~\ref{lem:powers},
\[
\delta^s\xi=\Bigl(\sum_kd_kP_{k+s}\Bigr)\delta+\sum_kd_kP_{k+s-1},
\]
so the coefficient $p=\sum_kd_kP_{k+s}$ of $\delta$ depends only on $\xi$ (the coordinates in the basis
$\{\delta,1\}$ are unique), not on the string. The shift $k\mapsto k+s$ therefore maps the fibre over
$\xi$ injectively into the set of words $(d_2,\dots,d_{4m+2})$ with $\sum_jd_jP_j=p$, which has
$\rho_{4m+1}(p)\le B_{4m+1}=4^m$ elements by \PETone. The bound is attained at the integer target
$C_m$ for $m\ge1$ and at $0,1,2$ for $m=0$, by \STTtwo.
\end{proof}

\begin{remark}[scope of \STTthree]\label{rem:scope11}
\STTthree{} concerns all targets in $\Z[\sqrt2]$, not only integers. It is stronger than what
\STTone{} gives directly, since \STTone{} says nothing about $\xi\in\Z[\sqrt2]\setminus\Z$. No target
outside $\Z[\sqrt2]$ is represented by any string, so the fibre there is empty. The bound concerns
the stated alphabet and window only.
\end{remark}

%======================================================================
\section{Discussion and limitations}\label{sec:discussion}
%======================================================================

\subsection{Scope}
\PETone{} and \PETtwo{} treat all lengths and all odd maximizers; \PETthree{} and \PETfour{} treat the
even maximizers, with the exceptional base case $k=1$. Theorems STT-1 to STT-3 concern the specified
ternary window and the target domains stated in the theorems: rational integers in \STTone{} and
\STTtwo, and $\Z[\sqrt2]$ in \STTthree. The case $m=0$ of \STTtwo{} is separate.

\subsection{What is not claimed}
We give no description of $\rho_\ell(n)$ for general $n$, no finite automaton computing it, and no
asymptotic distribution of the coefficients. We do not describe coefficients close to the maximum.
We do not claim that the parametrization of \PETfour{} is optimal or minimal. We do not extend any
result to other recurrences, other digit alphabets, other windows or other Pisot units. We do not
claim that any statement in this paper is new.

\subsection{Prior art and the depth of reading}\label{sec:prior}
The \emph{type} of question is not new: the maximum of the representation function at a fixed
numeration length, together with a description of the integers attaining it, is treated for the
Fibonacci system in~\cite{KMP2005} and for the $m$-bonacci systems in~\cite{KMP2007}, and an exact
maximum for minimal-weight Fibonacci expansions is~\cite[Theorem~4.6]{GrabnerSteiner2011}. In the
sources listed below we did not find a treatment of the Pell system with digit set $\{0,1,2\}$ on the
bounded window $P_2,\dots,P_{\ell+1}$, nor of fibres in $\Z[\sqrt2]$ over a symmetric window of
exponents. No count-preserving transfer from the Fibonacci or $m$-bonacci results to the present
setting is known to us: the recurrences ($x^2-x-1$ against $x^2-2x-1$), the alphabets and the
normalizations differ.

\begin{center}
\small
\begin{tabularx}{\textwidth}{@{}l>{\raggedright\arraybackslash}X>{\raggedright\arraybackslash}X@{}}
\toprule
Source & Depth of reading in the project & Relation to this paper\\
\midrule
\cite{Horadam1994} & full text & one distinguished (lazy) representation; not a fibre count\\
\cite{Horadam1996} & zbMATH review and metadata only & distinguished min/max Pell representations\\
\cite{KnopfmacherRobbins2004} & full text & greedy representation; partitions of unbounded multiplicity\\
\cite{KMP2005} & full text & Fibonacci, digits $\{0,1\}$: maximum and maximizers; closest type of result\\
\cite{KMP2007} & full text & $m$-bonacci, digits $\{0,1\}$: same type of question\\
\cite{Stockmeyer2009} & full text & sharp bound $R(n)\le\sqrt{n+1}$\\
\cite{GrabnerSteiner2011} & full text & minimal-weight expansions only\\
\cite{Edson2013} & publisher abstract only & general linear numeration systems; transfer not established\\
\bottomrule
\end{tabularx}
\end{center}

The full texts were read at earlier stages of the project; later stages relied on that reading and
on bibliographic records, and the present text was edited without rereading them. A structured
search of zbMATH Open, Crossref, OpenAlex, arXiv and OEIS was made at one stage. The remaining gaps
are as follows. MathSciNet, Web of Science and Scopus were not accessed, and Google Scholar and
Semantic Scholar were not usable. The full texts of~\cite{Horadam1996} and~\cite{Edson2013} were
not available. Two 2026 preprints in a related
setting were located (one on a Pell-to-ternary numeration map, read at the level of its description
only, and one on a silver-ratio sieve, not retrieved), and a 2019 paper on ternary strings and Pell
numbers was identified by title only; any of these, or a work not found, could contain an equivalent
statement. A human bibliographic check is required before any claim of novelty.

\subsection{Verification status}\label{sec:verification}
Every proof in this paper is self-contained and uses no computation. The proofs have been examined
by AI models only (see the statement on AI contribution below); they have not been read by a human
mathematician, and no formal (machine-checked) proof exists. Exact computations made at earlier
stages of the project---coefficients of $G_\ell$ for $\ell\le16$ by two independent methods, the
silver-ratio fibres computed directly in $\Z[\sqrt2]$ without using \STTone, the set recursions of
Sections~\ref{sec:pet3}--\ref{sec:pet4} for large $k$, and the identities of
Sections~\ref{sec:stt1}--\ref{sec:stt2} for $m\le300$---agree with every statement. During the
preparation of this text the formulas, indices and worked examples were checked again against exact
computation: $\rho_\ell$ for $\ell\le17$, $\Fam_k=M_{2k}$ for $k\le8$, the set recursions of
Lemma~\ref{lem:Y} and Proposition~\ref{prop:Yclosed} for $k\le200$, and the fibres $\nu_m$ directly in
$\Z[\sqrt2]$ for $m\le4$. These computations are falsification attempts and reproducibility aids,
not proofs, and nothing in Sections~\ref{sec:pet1}--\ref{sec:stt3} depends on them.

\subsection{Open questions}
\begin{enumerate}[label=(\arabic*)]
\item Describe the coefficients of $G_\ell$ just below the maximum, for instance the second largest
value and its arguments.
\item Is there a simpler description of $M_{2k}$, for example one without the partial sums $S_j$?
\item Determine for which windows, alphabets and recurrences the window lemma
(Lemma~\ref{lem:window}) keeps a positive margin, and what replaces it otherwise.
\item Which elements of $\Z[\sqrt2]\setminus\Z$ attain the bound of \STTthree?
\end{enumerate}

%======================================================================
\section*{Statement on AI contribution and provenance}
%======================================================================

The proofs and the principal results of this paper were developed with AI systems, in a staged
research project directed by a human project owner.
The following attribution is taken from the project's controlled records; where those records do not
establish an executor, we say so.
\begin{itemize}[leftmargin=1.6em]
\item \emph{Early stages} (the silver-ratio model and the first derivation of the transfer underlying
\STTone): the executing system is not recorded in the stage reports (\textsc{not established}).
\item \emph{Stage P7}: OpenAI \texttt{gpt-6-luna} in the Codex environment derived \PETone{} and
\PETtwo{} for $\ell\ge3$ and the silver-ratio maximum and uniqueness argument used in \STTtwo.
\item \emph{Stage P8}: Anthropic \texttt{claude-opus-5-5} in Claude Code carried out a model
certification of P7, extended \PETtwo{} to $\ell=1$, gave the proof of \STTone{} used here (by a
route different from that of the early stages) and derived \STTthree.
\item \emph{Stages P9--P10}: DeepSeek \texttt{deepseek-flash} in the DeepSeek Harness carried out a
prior-art audit and manuscript revision (P9); the executor of the bibliographic follow-up and
full-text reading (P10) is not established in the records.
\item \emph{Stage P11}: OpenAI \texttt{gpt-6-luna} in Codex proved the even-length set recursion and
\PETthree.
\item \emph{Stage P12}: Anthropic \texttt{claude-opus-5-5} in Claude Code audited P11 and proved the
explicit parametrization \PETfour.
\item \emph{Stage P13}: DeepSeek \texttt{deepseek-flash} in the DeepSeek Harness independently
certified \PETfour{} at the level of the executing model: it is a different provider, model and
harness from the author of P12. This certification was not blind (the proof was read before it was
reconstructed) and it is not a human review.
\item \emph{Stage P14}: a consolidation stage, recorded in its own register as OpenAI
\texttt{gpt-6.1-terra} in Codex, reconciled the theorem labels, adopted documentary corrections
(numerical, indexing and wording errata that changed no statement or proof) and assembled the working
manuscript from which this paper was edited. It was not a new mathematical audit.
\item \emph{This text} was edited from that working manuscript by Claude (Anthropic), in a session
configured for the model identifier \texttt{claude-opus-5-5}, the same model identifier as the
executor of P8 and P12. The editing is not a new independent certification of the proofs. It
corrected two index errors in side remarks (in Section~\ref{sec:pet1} and
Remark~\ref{rem:special}) and two bibliographic details (a page range and a misattributed review
number); none of these affects a statement or a proof.
\end{itemize}
Agreement between AI systems is not human peer review, and executor-level independence does not
imply independence of training. No human has reviewed the proofs, no human proof-verification or
citation-check certificate exists, and no priority is claimed. No AI system is presented as a human
author.

\bibliographystyle{amsalpha}
\bibliography{HGS-SILVER-REFERENCES}

\appendix

%======================================================================
\section{Worked small cases}\label{app:examples}
%======================================================================

\subsection*{\texorpdfstring{A.1. $\ell=1,2$}{A.1. l=1,2}}
$G_1=1+x^2+x^4$, $B_1=1$, $M_1=\{0,2,4\}$. $G_2=(1+x^2+x^4)(1+x^5+x^{10})$, $B_2=1$, with the nine
distinct exponents $0,2,4,5,7,9,10,12,14$, so $M_2$ has $9$ elements.

\subsection*{\texorpdfstring{A.2. $\ell=3$}{A.2. l=3}}
Weights $2,5,12$; $D_3=38$, $B_3=2$. By \PETtwo{} with $u_3=12$ and $U=P_4=12$,
$M_3=\{12,14,24,26\}$. The coefficients $\rho_3(0),\dots,\rho_3(38)$ are
\[
1,0,1,0,1,1,0,1,0,1,1,0,2,0,2,0,1,1,0,1,0,1,1,0,2,0,2,0,1,1,0,1,0,1,1,0,1,0,1 .
\]

\subsection*{\texorpdfstring{A.3. $\ell=4$ ($k=2$)}{A.3. l=4 (k=2)}}
Here $a=P_6=70$, $c=P_5=29$, $z_1=29$, $z_2=58$, $S_3=5$, $S_4=17$, $D_4=96$, $B_4=2$. \PETfour{}
gives
\[
\begin{aligned}
&\OL_4=\{12,14\},\ \OL_3=\{24,26\},\ \CC_1=\{29,31,33,34,36,38\},\ \IR_3=\{41,43\},\\
&\IL_3=\{53,55\},\ \CC_2=\{58,60,62,63,65,67\},\ \OR_3=\{70,72\},\ \OR_4=\{82,84\},
\end{aligned}
\]
$24=16\cdot2-8$ maximizers in all, symmetric about $D_4/2=48$, with $\min M_4=v=12$.

\subsection*{\texorpdfstring{A.4. $m=1$ ($\ell=5$, $s=4$, $P_s=12$)}{A.4. m=1 (l=5, s=4, P\_s=12)}}
$u_5=(P_7-5)/2=82$ and $U=P_6=70$, so $M_5=\{82,84,152,154\}$. Further $H_1=D_5=P_7+P_6-3=236\equiv-4
\pmod{12}$. The only maximizer divisible by $12$ is $84=7\cdot12$, so $C_1=7$ and $\nu_1(7)=4$.

\subsection*{\texorpdfstring{A.5. Small $m$}{A.5. Small m}}
\[
\begin{array}{c|c|c|c|c}
m & s & P_s & C_m & N_m\\\hline
0 & 2 & 2 & 1 & 1\\
1 & 4 & 12 & 7 & 4\\
2 & 6 & 70 & 41 & 16\\
3 & 8 & 408 & 239 & 64\\
4 & 10 & 2378 & 1393 & 256
\end{array}
\]
For $m=0$ the value $C_0=1$ is one of the three maximizers $0,1,2$.

%======================================================================
\section{What is and is not established}\label{app:status}
%======================================================================

\begin{center}
\small
\begin{tabularx}{\textwidth}{@{}>{\raggedright\arraybackslash}X l l l@{}}
\toprule
Statement & Established & Range & Where\\
\midrule
$B_\ell=2^{\floor{(\ell-1)/2}}$ and the pair invariant~\eqref{eq:pair} & yes & all $\ell\ge1$ & \PETone\\
complete list of maximizers & yes & odd $\ell\ge1$ & \PETtwo\\
number of maximizers, $9$ and $16k-8$ & yes & $k=1$; $k\ge2$ & \PETthree\\
explicit description of $M_{2k}$ & yes & $k\ge1$ & \PETfour\\
disjointness of the six families & yes & $k\ge2$ (fails at $k=1$) & Prop.~\ref{prop:order}\\
$\nu_m(t)=\rho_{4m+1}(tP_{2m+2})$ & yes & $m\ge0$, $t\in\Z$ & \STTone\\
$\max_t\nu_m(t)=4^m$ & yes & $m\ge0$ & \STTtwo\\
unique maximizing integer $C_m$ & yes & $m\ge1$ & \STTtwo\\
$\argmax\nu_0=\{0,1,2\}$ & yes & $m=0$ & \STTtwo\\
every $\Z[\sqrt2]$-fibre has at most $4^m$ strings & yes & $m\ge0$ & \STTthree\\
optimality of the parametrization & no & --- & not claimed\\
coefficients near the maximum, automata, asymptotics & no & --- & not attempted\\
other recurrences, alphabets, windows & no & --- & not attempted\\
priority or novelty & no & --- & unresolved, Section~\ref{sec:prior}\\
human review & no & --- & not requested\\
\bottomrule
\end{tabularx}
\end{center}

%======================================================================
\section{Canonical identifiers, historical labels and computations}\label{app:aliases}
%======================================================================

The project records used other labels for some theorems. To avoid ambiguity, the canonical
identifiers of this paper correspond to them as follows; in particular the historical
``Theorem~C'' and ``Theorem~D'' always denote the silver-ratio transfer and the silver-ratio maximum,
never the even-length results.

\begin{center}
\small
\begin{tabular}{@{}lll@{}}
\toprule
Canonical identifier & Content & Historical label\\
\midrule
PET-1 & exact maximum coefficient & Theorem A\\
PET-2 & odd-length maximizers & Theorem B\\
PET-3 & number of even-length maximizers & even-length count (stage P11)\\
PET-4 & explicit even-length parametrization & P12-B (``B-LEVEL-3'')\\
STT-1 & silver-to-Pell transfer & Theorem C\\
STT-2 & silver maximum and uniqueness & Theorem D\\
STT-3 & algebraic fibre bound & Corollary E\\
\bottomrule
\end{tabular}
\end{center}

\emph{Role of computations.} The project's computer checks recompute, with exact integer arithmetic,
the degree formula and the gap pattern of Lemma~\ref{lem:shifts}; $B_\ell$ and $M_\ell$ for small
$\ell$; the two-step recursion~\eqref{eq:twostep} coefficient by coefficient; the silver-ratio fibres
directly in $\Z[\sqrt2]$; the quantities $C_m$, $z_m$, $H_m$ and the residues of
Lemma~\ref{lem:residue}; and the families of \PETfour. They cannot establish any statement with a free
parameter. Their functions are to let a reader falsify a claim cheaply and to make the computations
cited in Section~\ref{sec:verification} reproducible. A reader who finds a discrepancy between a
computation and the text should trust neither and re-derive the step by hand.

\end{document}
